\documentclass[10pt]{article}

\usepackage[utf8]{inputenc}

\usepackage[T1]{fontenc}

\usepackage{amsmath}

\usepackage{amsfonts}

\usepackage{amssymb}

\usepackage[version=4]{mhchem}

\usepackage{stmaryrd}

\usepackage{fixltx2e}



\DeclareUnicodeCharacter{0131}{$\imath$}



\begin{document}

Элемент площади:



\[

d \sigma=4\left(u^{2}+v^{2}\right) d u d v

\]



Векторные диффферевциальные онерации:



\[

\begin{aligned}

& \operatorname{grad}_{u} f=\frac{1}{2 \sqrt{u^{2}+v^{2}}} \frac{\partial f}{\partial u}, \\

& \operatorname{grad}_{v} f=\frac{1}{2 \sqrt{u^{2}+v^{2}}} \frac{\partial f}{\partial v},

\end{aligned}

\]



$\operatorname{div} \boldsymbol{a}=\frac{1}{2 \sqrt{u^{2}+v^{2}}}\left(\frac{\partial a_{u}}{\partial u}+\frac{\partial a_{v}}{\partial v}\right)+\frac{1}{2 \sqrt{\left(u^{2}+v^{2}\right)^{i}}}\left(u a_{u}+v a_{v}\right)$,



\[

\Delta f=\frac{1}{l\left(u^{2}+v^{2}\right)}\left(\frac{\partial^{2} f}{\partial u^{2}}+\frac{\partial^{2} f}{\partial v^{2}}\right) \text {. }

\]



У равнение Лапласа допускает в П. к. разделение переменных.



ПАРАБОЛИЧЕСКИЙ ЦИЛИНДР - цилиндрическая поверхность второго порядка, для к-рой паправляющей служит парабола. Каноничское уравнение II. ц.:



\[

y^{2}=2 p x

\]



ПАРАБОЛИЧЕСКОГО ТИПА УРАВНЕНИЕ - У. УАавнение вида



\[

\begin{gathered}

u_{t}-\sum_{i, j=1}^{n} a_{i j}(x, t) u_{x_{i} x_{j}}-\sum_{i=1}^{n} a_{i}(r, t) u_{x_{i}}- \\

-a(x, t) u=f(x, t),

\end{gathered}

\]



где $\sum a_{i} \xi_{i} \xi_{j}$ - положительно определенная квадратичная форма. Переменная $t$ выделена и играет роль времени. Типичным примером П. т. у. является уравненис теплопроводности



\[

u_{t}-\sum_{i=1}^{n} u_{x_{i} x_{i}}=0

\]



А. П. Солдатов.



ПАРАБОЛИЧЕСКОГО ТИПА УРАВНЕНИЕ; ч И сле н ны е ме т о ды ре ше ни я- методы решения уравненшй параболич. тиша на основе вычислительных алгоритмов. Для решения П. т. у. часто применяются приближенные численные методы, рассчитанные на использование быстродействующих ЈВМ. Наиболее универсальным является ме тод с с то к (к о н е чн 0 p а 3 н о с т ны й м с т о д).



Ниже рассмотрен метод сеток на иримере уравнения тешлошроводности



\[

\frac{\partial u}{\partial t}=\frac{\partial^{2} u}{\partial x^{2}}+f(x, t), \quad t>0, \quad 0<x<l,

\]



с краевыми условиями 1-го рода



\[

u(0, t)=\mu_{1}(t), u(l, t)=\mu_{2}(t), u(x, 0)=u_{0}(x) .

\]



Вводится равномерная сетка узлов $\left(x_{i}, t_{n}\right)$,



$x_{i}-i h, i=0,1, \ldots, N, h N=l, t_{n}=n \tau, n=0,1, \ldots ;$



If обозначения



\[

\begin{gathered}

y_{i}^{n}=y\left(x_{i}, t_{n}\right), y_{t, i}=\left(y_{i}^{n+1}-y_{i}^{n}\right) / \tau, \\

y \frac{n}{x_{1} i}=\left(y_{i}^{n}-y_{i-1}^{n}\right) / h, y_{x, i}^{n}=\left(y_{i+1}^{n}-y_{i}^{n}\right) / h, \\

y \frac{n}{x x, i}=\left(y_{i+1}^{n}-2 y_{i}^{n}+y_{i-1}^{n}\right) / h^{2} .

\end{gathered}

\]



Метод сеток состоит в том, что уравнение (1) приближенно заменяется системой линейных алгебраич. уравнений (разностной схемой)



$y_{t, i}^{n}=\sigma y_{\bar{x} x, i}^{n+1}+(1-\sigma) y_{\bar{x} x, i}^{n}+\varphi_{i}^{n}, i=1,2, \ldots, N-1$,

$y_{0}^{n}=\mu_{1}\left(t_{n}\right), y_{N}^{n}=\mu_{2}\left(t_{n}\right), y_{i}^{0}=u_{0}\left(x_{i}\right)$,



гцс $\sigma$ - числової параметр й $\varphi_{i}^{n}$ - сеточная аıицоксимация фұункции $f(x, t)$, напр. $\varphi_{i}^{n}=f\left(x_{i}, t_{n}-1-0,5 \tau\right)$. Система уравнений (2) репается по слоям, т. е. іля каждого $n=0,1, \ldots$ по известным значениям $y_{i}^{n}$, œ $i$ находятся новыс значения $y_{i}^{n+1}, i=1,2, \ldots, N-1$. Если $\sigma=0$ (лвная схема), то $y_{i}^{n+1}$ выражкаются явным образом через $y_{i}^{n}, \varphi_{i}^{n}$. Если же $\sigma \neq 0$ (неявныс схемы), то относительно $y_{i}^{n+1}, i=1,2, \ldots, N-1$, возникает система уравнений, имеющая трехдиагональную матрицу. Эта система уравнений решается прогонки методом. Недостатком явной схемы является сильное ограничение на лшаг $\tau$, возникающее из условия устойчивости, а именно $\tau \leqslant 0,5 h^{2}$. Неявные схемы при $\sigma \geqslant 0,5$ абсолютно устойчивы, т. е. устойчивы при любых шагах $h$ и $\tau$. Известны и другие разностные схемы для уравнения (1) (см. [1], [2]).



Если схема (2) устойчива и $\varphi_{i}^{n}$ ашпроксимирует $f(x, t)$, то при $h, \tau \rightarrow 0$ решение $y_{i}^{n}$ разностной задачи сходится к решению $u\left(x_{i}, t_{n}\right)$ исходной задачи (см. [1]). Порядок точности зависит от параметра $\sigma$. Так, симметритная схема $\left(\sigma=0,5, \varphi_{i}^{n} \cdot f\left(x_{i}, t_{n}+0,5 \tau\right)\right)$ имеет второй порядок точности по $\tau$ іl по $h$, т. е. для гаждого $n$ выполняется оценка



\[

\max _{0<i<\mathrm{N}}\left|y_{i}^{n}-u\left(x_{i}, t_{n}\right)\right| \leqslant M\left(\tau^{2}+h^{2}\right)

\]



с константой $M$, не зависнщей от $h$ и $\tau$. II ри $\sigma=0,5-$ $-h^{2} /(12 \tau)$ и специальном выборе $\varphi_{i}^{n}$ схема (2) имеет точность $O\left(\tau^{2}+h^{4}\right)$. Для остальных $\sigma-$ тотность $O\left(\tau+h^{2}\right)$



При решении П. т. у. на больших отрезках времени существенное значение имеет асимптотич. устойчивость разностной схемы. Решение уравнения (1) с $f(x, t)=0$, $\mu_{1}(t)=\mu_{2}(t)=0$ ведет себя при $t \rightarrow \infty$ как $e^{-\lambda_{1} t}$, $\lambda_{1}=\pi^{2} / l^{2}$. Этим свойством обладает не всякая разностная схема, аппроксимирующая уравнение (1). Напр., симметричная схема $(\sigma=0,5)$ асимптотически устойчива при условии $\tau \leqslant l h / \pi$. Построены схемы 2-го порядка точности, асимптотически устойчивые при любых $\tau$, $h$, однако они не входят в семейство (2) (см. [3]).



Краевые условия 3-го рода



\[

\frac{\partial u}{\partial x}(0, t)-\beta u(0, t)=\mu(t)

\]



апороксимируются с порядком $O\left(h^{2}\right)$ разностными у равнениями



$\sigma\left(y_{x, 0}^{n+1}-\left(\beta y_{0}^{n+1}+\mu^{n+1}\right)\right)+(1-\sigma)\left(y_{x, 0}^{n}-\left(\beta y_{0}^{n}+\mu^{n}\right)\right)=$



\[

=0,5 h\left(y_{t, 0}^{n}-\varphi_{0}^{n}\right), \mu^{n}=\mu\left(t_{n}\right) \text {. }

\]



Для уравнсния теплопроводности в цилиндрич. и сферич. координатах



\[

\frac{\partial u}{\partial t}=\frac{1}{r^{m}} \frac{\partial}{\partial r}\left(r^{m} \frac{\partial u}{\partial r}\right), m=1,2,0<r<R,

\]



вводятся сетки $\left(r_{i}, t_{n}\right)$, где $t_{n}=n \tau, n=0,1, \ldots$,



$r_{i}=\left(i+\frac{m}{2}\right) h, i=0,1, \ldots, N, h=\frac{R}{N+0,5 m}, m=1,2$, и рассматриваются разностные уравнения (см. [4])



\[

y_{t, i}^{n}=\Lambda_{r}^{(m)}\left(\sigma y_{i}^{n+1}+(1-\sigma) y_{i}^{n}\right)

\]



г де



\[

\Lambda_{r}^{(m)} y_{0}^{n}=\frac{1}{r_{0}^{m}}\left(\bar{r}_{1}^{m} \frac{y_{1}^{n}-y_{0}^{n}}{h}\right)

\][



\textbackslash Lambda\_\{r\}\^{}\{(m)\} y\_\{i\}\textsuperscript{\{n\}=\textbackslash frac\{1\}\{r\_\{i\}}\{m\}\} \textbackslash frac\{1\}\{h\}\textbackslash left(\textbackslash bar\{r\}\textit{\{i+1\}\^{}\{m\} \textbackslash frac\{y}\{i+1\}\textsuperscript{\{n\}-y\_\{i\}}\{n\}\}\{h\}-\textbackslash bar\{r\}\textit{\{i\}\^{}\{m\} \textbackslash frac\{y}\{i\}\textsuperscript{\{n\}-y\_\{i-1\}}\{n\}\}\{h\}\textbackslash right),

]



$\overline{r_{i}}=0,5\left(r_{i}+r_{i+1}\right)$ при $m=1, \bar{r}_{i}=\sqrt{r_{i} r_{i-1}} \quad$ при $m=2$, $i=1,2, \ldots, N-1$.





\end{document}